\section{Introduction}
\label{sec:introduction}

Law enforcement agencies (LEAs) rely on digital evidence, and many investigations now cross national borders \citep{hargreaves2024dfpulse}. Legal instruments already allow agencies to exchange evidence across borders \citep{eio2014,eevidence2023, reviewOfreview}. Training a machine learning model together is a different matter: a model trained on data from several agencies may reveal patterns that no single agency can see, but the training data is sensitive and often cannot be collected in one place \citep{scanlon2023chatgpt,breitinger2023sharing}. Federated Learning (FL) addresses this technically, since each agency keeps its data and only model updates are shared \citep{mcmahan2017communication}. However, existing FL frameworks were not built for forensic work: auditability, governance across organizational levels, explainability \citep{solanke2022explainable}, and evidentiary accountability are rarely design goals \citep{riedel2024comparative}.

Some legal requirements can be met through system design, for example by minimizing data, logging processing, and restricting access. Others, such as the legal basis, purpose limitation, and lawful further use, require legal assessment \citep{directive2016680,TeMolder2023}. Cross-border cooperation adds a specific question. In FL, evidence is always analyzed locally, so the question is whether information derived from it may contribute to a model that other jurisdictions receive.

Here, a flat FL topology reaches a limit. It can comply with a restriction by leaving the restricted material out of training. But it produces only one shared model, so the material is then also missing from the model used inside the jurisdiction, where its use is lawful and often most valuable. A flat topology can exclude restricted material, but it cannot scope where its influence travels. Scoping requires an aggregation point at the border of the jurisdiction, which is a property of the topology, not a setting \citep{9148862}.

This paper therefore asks: \emph{which federated learning topologies allow a jurisdictional restriction to be enforced at the point where it legally applies, and what does enforcing it cost in model quality?}

We make three contributions. First, we derive seven requirements (R1 to R7) for FL in forensic settings, state what the forensic context adds beyond other regulated domains, identify a conflict between two of them, and assess the open-source FL ecosystem against them. Second, we propose \emph{Federated Forensics}, a three-layer reference architecture whose middle layer sits at the border of legal competence and can withhold a contribution from further aggregation. Third, we evaluate a prototype in a controlled experiment: the gate blocked every restricted update, and accuracy on the restricted classes rose from 0\% under the compliant flat alternative to 93.0\%, 4.4 points below ignoring the restriction.

Two delimitations apply. First, our results are intelligence, not evidence: a model update loses its link to specific items of raw evidence once it leaves the agency, so a finding can direct evidence collection under chain-of-custody procedures but never replaces it. Second, we do not design the audit mechanism. The architecture fixes where audit records must be written, but their format, verification, and conflict with differential privacy (Sec.~\ref{subsec:req_tensions}) are treated in separate work.

Sec.~\ref{Federated Learning and Digital Forensics: Context and Gaps} positions the work, Sec.~\ref{Forensic Specific Requirements for Federated Learning} derives the requirements, Sec.~\ref{Federated Forensics: A Reference Architecture for Law Enforcement Collaboration} presents the architecture, Sec.~\ref{Prototype Implementation} evaluates it, and Secs.~\ref{sec:limitations} and~\ref{sec:conclusion} discuss limitations and conclude.
